
\lesson{13/10}
In this formalism, the equilibrium is found when the explicit dependency of $\rho$
on time is null:
$$\pdv[]{\rho}{t}=0 \leftrightarrow \left\{\rho,H\right\}=0$$
which means that $\rho$ should be a \textit{constant of motion}, depending on
$q,p$ and other constants of motion.

\section{Description of TD systems through independent subsystems}
Consider a TD system composed of quasi-isolated subsystems (i.e. the interaction
between subsystems are confined to their boundaries): for this reason, their
probability densities can be well approximated as independent, which means
$$\rho_{ij}(q_i,p_i,q_j,p_j) = \rho_i(q_i,p_i)\rho_j(q_j,p_j)$$
for every subsystems $i\neq j$.
\subsection{Additive observables}
These observables can be expressed as the sum of the observable in each subsystem
$$O = \sum_{i=1}^{n}O_i,$$
like the total internal energy; their expectation values are also additive.

For this approach to be useful, the inevitable fluctuations about the equilibrium
state must be small: consider the variance
\begin{gather*}
\sigma^2 = \langle O^2\rangle-\langle O\rangle^2 = \langle\sum_{ij}^{}O_iO_j
\rangle - \langle\sum_{i}^{}O_i\rangle^2 = \sum_{ij}^{}\left(
\langle O_iO_j\rangle - \langle O_i\rangle\langle O_j\rangle\right) =\\=
\sum_{i}^{}\left(\langle O_i^2\rangle-\langle O_i\rangle^2\right) +
\sum_{i\neq j}^{}\left(\langle O_iO_j\rangle - \langle O_i\rangle\langle O_j\rangle
\right) = \sum_{i}^{}\sigma^2_i
\end{gather*}
where the last sum vanishes for the quasi-isolation condition.

The mean value scales with $n$, while the standard deviation scales like
$\sqrt{n}$: their quotient becomes smaller as $n$ gets bigger, so the statistical
approach becomes better the bigger the number of subsystems.
\subsection{Explicit formula for $\rho$}
Since $\rho_{ij} = \rho_i\rho_j$, $\ln(\rho_{ij}) = \ln(\rho_i)+\ln(\rho_j)$ is
an extensive quantity and so should only depend on \textit{additive} constants
of motion.

From classical mechanics it is known that there exist \textbf{only 7 additive
constants of motion}:
$$E_i,\,\mathbf{P}_i,\,\mathbf{L}_i$$
for the $i$'th subsystem; so the log of the probability density must be a
function in the form
$$\ln(\rho_i) = \alpha+\beta E_i + \mathbf{\gamma}\cdot\mathbf{P_i}+
\mathbf{\delta}\cdot\mathbf{L_i}$$
where the constants are independent of the subsystem and completely determined by
the \textit{total motion of the system}: the only non-trivial dependency of
$\rho$ at equilibrium is \textbf{energy}, because the total momentum and angular
momentum can be factored out with a change of reference frame.
\section{The Time reversal problem}
One of the main problems in statistical mechanics is the fact that microscopic
laws are time-reversible, while the macroscopic laws of TD are not.
\subsection{Reduced Observables}
A one-body observable is an observable in the form
$$O_1(q,p) = \sum_{i=1}^{N} O(\mathbf{q_i,p_i})$$
for example $E_k = \sum_{}^{}\flatfrac{\mathbf{p_i}^2}{2m}$.

A two-body observable is an observable in the form
$$O_2(q,p) = \sum_{i<j=1}^{N}O(\mathbf{q_i},\mathbf{p_i},\mathbf{q_j},
\mathbf{p_j}).$$

Consider the expected value of a 1-body observable
$$\langle O_1(t)\rangle = \int_{}^{}\dd\mu_1...\dd\mu_N \sum_{i=1}^{N} O(i)
\rho(1,2,...N)$$
the density does not distinguish the particles, so $\rho$ must be symmetric
for variable swap: this means that the sum can be factored out as a $N$ factor
$$\langle O_1(t)\rangle = N\int_{}^{}\dd\mu_1 O(1)
\underbrace{\int_{}^{}\dd\mu_2...\dd\mu_N \rho(1,2,...N)}_{:=\rho_1(1;t)}$$
this $\rho_1$ is called the \textbf{reduced probability density of order 1} and
it is clearly normalized.

The same can be done for two-body observables
$$\langle O_2(t)\rangle=\int_{}^{}\dd\mu_1...\dd\mu_N \sum_{i<j}^{}O(i,j)\rho
(1,2...N)$$
The double sum is comprised of $\begin{pmatrix}
N\\2
\end{pmatrix}
$ equal terms, so the \textbf{reduced probability of order 2} is 
$$\begin{pmatrix}
N\\2
\end{pmatrix}
\int_{}^{}\dd\mu_1\dd\mu_2 O(1,2)\underbrace{\int_{}^{}\dd\mu_3...\dd\mu_N
\rho(1,2,...N)}_{:=\rho_2(1,2;t)}
$$

\textbf{Assumptions} (Diluted Gas):
\begin{itemize}
    \item $H = \sum_{i}^{}\frac{\mathbf{p}_i^2}{2m} + \sum_{i<j}^{}\phi_{ij}$;
    \item $\phi_{ij} = \phi_{ji} = \phi(\left|\mathbf{q}_i-\mathbf{q_j}\right|)$
    \item "Short range" potential: $\phi_{ij} = 0$ for $\left|\mathbf{q}_i-
    \mathbf{q_j}\right|>r_0$
\end{itemize}
The Hamilton equations are
$$\begin{dcases*}
  \pdv{H}{\mathbf{p}_i} = \frac{\mathbf{p}_i}{m}\\
  \pdv{H}{\mathbf{q}_i} = \sum_{k<l}^{}\pdv{\phi_{kl}}{\mathbf{q}_i}
  \end{dcases*}$$
the second equation can be simplified to
$$\pdv{H}{\mathbf{q}_i} = \sum_{i<l}^{}\pdv{\phi_{il}}{\mathbf{q}_i} + 
\underbrace{\sum_{k<i}^{}\pdv{\phi_{ki}}{\mathbf{q}_i}}_{=\sum_{i<k}^{}...} = 
\sum_{j\neq i}^{} \pdv{\phi_{ij}}{\mathbf{q}_i} := -\mathbf{K}_{ij}$$
where $\mathbf{K}_{ij} = \mathbf{K}_{ji}$: Liouville's theorem becomes
$$\pdv{\rho}{t} + \sum_{i}^{}\frac{\mathbf{p}_i}{m}\pdv{\rho}{\mathbf{q}_i} + 
\sum_{i\neq j}^{}\mathbf{K_{ij}}\pdv{\rho}{\mathbf{p}_i}=0$$
antisymmetrizing the second sum yields
$$\pdv{\rho}{t} + \sum_{i}^{}\frac{\mathbf{p}_i}{m}\pdv{\rho}{\mathbf{q}_i} + 
\frac{1}{2}\sum_{i\neq j}^{}\mathbf{K_{ij}}\left(\pdv{\rho}{\mathbf{p}_i}
\pdv{\rho}{\mathbf{p}_j}\right)=0.$$
Defining the operators
\begin{gather*}
\hat S_i = \frac{\mathbf{p}_i}{m}\pdv[]{}{\mathbf{q}_i}\\
\hat P_{ij} = K_{ij}\left(\pdv{}{\mathbf{p}_i}-\pdv[]{}{\mathbf{p}_j}\right)\\
\end{gather*}
the equation becomes
$$(\pdv[]{}{t} + \underbrace{\sum_{i=1}^{N}\hat S_i +
\frac{1}{2}\sum_{i\neq j=1}^{N}\hat P_{ij}}_{:=\hat L_N})\rho=0.$$
To find the consequences of this calculation on the reduced probability density,
it is useful to redefine them with a different normalization:
\begin{gather*}
f_1(i) = N\rho_1(i)\\
f_s(1,2,...s) = \frac{N!}{(N-S)!}\int_{}^{}\dd\mu_{s+1}...\dd\mu_N\rho(1,2...N)=
\frac{N!}{(N-S)!}\rho_s(1,2...s)
\end{gather*}
where the second formula is the generalization of the reduced probability density
at $s$ bodies.

The relation with the Liouville operator $\hat L_N$ is
$$\pdv[]{f_s}{t} = -\frac{N!}{(N-S)!}\int\dd\mu_{s+1}...\dd\mu_N\hat L_N
\rho$$
It is better to subdivide $\hat L_N$ in the following fashion
$$
\hat L_N = \sum_{i=1}^{s}\hat S_i + \sum_{j=s+1}^{N}\hat S_j + \frac{1}{2}
\sum_{i\neq j=1}^{s}\hat P_{ij} + \frac{1}{2}\sum_{i\neq j=s+1}^{N}\hat P_{ij} +
\sum_{i=1}^{s}\sum_{j=s+1}^{N} \hat P_{ij}=\hat L_s + \hat L_{N-s}+\sum_{i=1}^{s}
\sum_{j=s+1}^{N}\hat P_{ij}
$$
So the derivative becomes
$$\pdv[]{f_s}{t}\frac{N!}{(N-s)!}\int_{}^{}\dd\mu_{s+1}...\dd\mu_N\left(\hat
L_s\rho + \underbrace{\hat L_{N-s}\rho}+\sum_{i=1}^{s}\sum_{j=s+1}^{N}
\hat P_{ij}\rho\right)$$
the highlighted part is zero, because its parts can be factored and integrated
by parts, e.g.
$$\int_{}^{}\dd^3p_{s+1}\dd^3q_{s+1} \frac{\mathbf{p}_{s+1}}{m}
\pdv{\rho}{\mathbf{q}_{s+1}}= [...?!]$$
and so reshuffling the terms
$$\left(\pdv[]{}{t}+\hat L_s\right)f_s = -\frac{N!}{(N-s)!}
\int_{}^{}\dd\mu_{s+1}...\dd\mu_N \sum_{i=1}^{s}\sum_{j=s+1}^{N}\hat P_{ij}\rho$$
since the $s+1,...,N$ coordinates are integrated, the sum over $j$ gives the same
result $N-s$ times, using $j=s+1$ yields
$$\left(\pdv[]{}{t}+\hat L_s\right)f_s = -\frac{N!}{(N-s)!}(N-s)\sum_{i=1}^{s}
\int_{}^{}\dd\mu_{s+1}...\dd\mu_N\hat P_{i,s+1} \rho=$$
$$= -\sum_{i=1}^{s}\int_{}^{}
\dd\mu_{s+1}\hat P_{i,s+1}\underbrace{\frac{N!}{(N-(s+1))!}\int_{}^{}
\dd\mu_{s+2}...\dd\mu_N\rho}_{:=f_{s+1}}=$$
$$=-\sum_{i=1}^{s}\int_{}^{}\dd\mu_{s+1}\mathbf{K}_{i,s+1}\left(
\pdv[]{}{\mathbf{p}_i}-\pdv[]{}{\mathbf{p}_{s+1}}\right)f_{s+1}(1,2...,s+1;t)$$
the derivative with respect to $\mathbf{p}_{s+1}$ can be cancelled integating
by parts.

The final result is the BBGKY hierarchy
$$\left(\pdv{}{t}+\hat L_s\right)f_s(1,2...s;t) = -\sum_{i=1}^{s}\int_{}^{}
\dd\mu_{s+1}\mathbf{K}_{i,s+1}\pdv[]{}{\mathbf{p}_i}f_{s+1}(1,2,...,s+1;t)$$
